% -*- TeX-master: "lectures-on-number-theory" -*-

\chapter{Quadratic Reciprocity}

\section{Quadratic residues}

A quadratic congruence is a congruence of the form
\begin{equation*}
  ax^2 + bx + c \equiv 0 \pmod{n}.
\end{equation*}
The simplest kind of quadratic congruence is of the form $x^2-c\equiv 0\pmod{n}$, or equivalently $x^2\equiv c\pmod{n}$. This congruence poses the question of establishing $c$ as the square of some other integer modulo $n$. For example, $0$ and $1$ are always squares modulo any natural number $n$. Modulo $3$ these are the only squares, since the quadratic congurence $x^2\equiv 2\pmod{3}$ does not have solutions. In other words, every square not divisible by $3$ must be congruent to $1$ modulo $3$, as can be readily verified by computing
\begin{equation*}
  (3k+j)^2=9k^2 + 6jk + j^2
\end{equation*}
for $1\leq j \leq 2$. Similarly, $0$ and $1$ are the only quadratic residues modulo $4$; that is, no square number can be congruent to $2$ or $3$ modulo $4$. Modulo $5$ we find that $0$, $1$, and $4$ are quadratic residues, and so on. The following table lists quadratic residues up to the number $8$:
\begin{center}
  \begin{tabular}{lllllllll}
    $n$ & $1$ & $2$ & $3$ & $4$ & $5$ & $6$ & $7$ & $8$ \\ \midrule
    quadratic residues & $0$ & $0,1$ & $0,1$ & $0,1$ & $0,1,4$ & $0,1,3,4$ & $0,1,2,4$ & $0,1,4$
  \end{tabular}
\end{center}
Consider an odd prime $p$. By \crefnamed{Lagrange's Theorem}{theorem:lagranges-theorem} it follows that the 


Euler had a deep understanding of numbers that are congruent to a square modulo a prime. In 1748 he had published a result now known as \emph{Euler's criterion}, which characterized precisely the squares modulo a prime.

\begin{theorem}[Euler's criterion]\label{theorem:eulers-criterion}
  Consider an integer $a$ and an odd prime $p$. Then we have
  \begin{equation*}
    a^{\frac{p-1}{2}} \equiv
    \begin{cases}
      1 & \text{if $x^2\equiv a\pmod{p}$ for some $x$} \\
      -1 & \text{otherwise.}
    \end{cases}
  \end{equation*}
\end{theorem}


\section{Euler's Prime Generating Polynomial}

In the early 1770s, Leonhard Euler divised his famous prime generating polynomial
\begin{equation*}
  n^2 + n + 41.
\end{equation*}
This polynomial produces prime numbers for any $0\leq n \leq 39$. Note that
\begin{equation*}
  40^2 + 40 + 41 = 40(40+1) + 41 = 41^2
\end{equation*}
isn't prime anymore. We could easily imagine Euler pondering one day over the question of creating a quadratic polynomial that returns a surprising amount of primes, and it would be natural to wonder how he might have approached the problem.

Of course, there is no polynomial $f(x)$ other than a constant polynomial such that $f(x)$ is a prime number for every integer $x$.

\begin{proposition}
  If every value of a polynomial $f(x)$ is prime, then $f(x)$ is constant.
\end{proposition}

\begin{proof}
  Suppose $f(x)$ is a polynomial, all of whose values are prime. Then in particular $p:=f(0)$ is prime. Then $p\mid f(kp)$ for every integer $k$. Since each value of $f$ is prime, it follows that $f(kp)=p$ for every integer $k$. This implies that $f(x)-p$ has infinitely many roots, i.e., $f(x)$ is the constant polynomial with value $p$.
\end{proof}

Thus we are faced with the task of finding a prime $q$ such that $n^2+n+q$ is prime for $0\leq n \leq q-2$. One thing to note is that if the number
\begin{equation*}
  n^2 + n+q
\end{equation*}
is composite for some $0\leq n\leq q-2$, then it must have a prime factor below $q$. Indeed, the integer $n^2+n+q$ is always below $q^2$, so any composite numbers below $q^2$ have a prime factor below $q$.

\begin{proposition}
  Consider a prime $p$. If there is a solution to the quadratic congruence
  \begin{equation*}
    ax^2 + bx + c \equiv 0 \pmod{p},
  \end{equation*}
  then the discriminant $\Delta=b^2-4ac$ is congruent to a square modulo $p$.
\end{proposition}

\begin{proof}
  Note that any number is a square modulo $2$, so we may assume that $p$ is an odd prime. Also note that if $p\mid a$, then $\Delta\equiv b^2\pmod{p}$, so we may assume that $p\nmid a$, i.e., that $a$ is invertible in $\mathbb{Z}/p\mathbb{Z}$. Since $p$ is odd, we also have that $4$ is invertible modulo $p$. The quadratic congruence in the statement of the proposition is therefore solvable if and only if the quadratic congruence
  \begin{equation*}
    4a^2x^2 + 4abx + 4ac \equiv 0 \pmod{p}
  \end{equation*}
  is solvable. Given a solution of this quadratic congruence, we calculate
  \begin{align*}
    (2ax + b)^2 & = 4a^2x^2 + 4abx + b^2 \\
                & = 4a^2x^2 + 4abx + 4ac + b^2 - 4ac \\
    & \equiv b^2 - 4ac
  \end{align*}
  modulo $p$.
\end{proof}

In other words, if the discriminant $\Delta$ of a polynomial $ax^2 + bx + c$ is not a square modulo $p$ for any prime $0\leq p<q$, then the polynomial congruence
\begin{equation*}
  ax^2 + bx + c \equiv 0 \pmod{p}
\end{equation*}
is impossible for any $0\leq p <q$.

This opens up a method of searching for a polynomial $n^2 + n + q$ returning primes for any $0\leq n \leq q-2$. The discriminant of this polynomial is $1-4q$, and our goal is to find a prime number $q$ such that $1-4q$ is not a square modulo $p$ for any $0\leq p < q$.

For example, the number $2$ is not a square modulo $3$, and the numbers $2$ and $3$ are both not squares modulo $5$.

\begin{exercises}
  \exercise Show that a number $n$ is square-free if and only if the only solution of the quadratic congruence $x^2\equiv 0\pmod{n}$ is $x=0$.
\end{exercises}
